\section{Conclusiones}\label{sec:conclusiones}

Con este trabajo se desarrollo CinemaBox, un asistente integrado con Telegram que permite a los usuarios realizar consultas y recibir recomendaciones de peliculas mediante dos metodos de recuperacion de informacion: TF-IDF y embeddings multilingues, ambos con la similitud del coseno. de esta manera, se logro cumplir con el objetivo principal del proyecto, que consistia en implementar una solucion que permitiera explorar un catalogo de peliculas mediante consultas realizadas.

Los datos se prepararon a partir de la informacion del dataset TMDB 5000 Movie Dataset, utilizando generos y sinopsis para representar el contenido de cada pelicula. Por otro lado, el metodo TF-IDF se implemento basandose en la importancia de los terminos y la similitud del coseno, los resultados obtenidos con palabras clave fueron ligeramente mejores que los de embeddings, aunque la diferencia no es concluyente según el bootstrap, mientras que el metodo de embeddings multilingues permitio representar las consultas y las peliculas en un espacio vectorial para comparar su similitud semantica, lo que permitio que presentara un comportamiento mas adecuado a la hora de recuperar peliculas a partir de las consultas en un lenguaje mas natural, como por ejemplo, descripciones de la trama de la pelicula. Por lo tanto, no se identifico un metodo que sea "superior" en todos los casos, sino que hay diferencias en su comportamiento dependieno de sus datos de entrada.

Los tiempos de respuesta fueron en promedio 35.1 ms para el metodo TF-IDF y 42.0 ms para embeddings, sin embargo, si se ejecuta embeddings con GPU, su tiempo promedio se reduce a 14.9 ms, lo que demuestra que el rendimiento del metodo esta condicionado por los recursos del hardware disponible, y que el uso de una GPU puede ayudar a disminuir el tiempo requerido para una respuesta.

Finalmente, la incorporación de un LLM permitió generar respuestas en lenguaje natural a partir de las películas recuperadas mediante embeddings. Aunque esto facilita una interacción más conversacional, las pruebas evidenciaron que el modelo puede seleccionar incorrectamente una película o añadir información que no está respaldada por las sinopsis proporcionadas.

En conclusion, el desarrollo del proyecto permitio implementar y comparar estos dos enfoques de recuperacion de informacion dentro de un asistente conversacional o bot, lo que permitio identificar las fortalezas y debilidades de cada metodo, siendo estas que por un lado TF-IDF resulta competitivo cuando las consultas comparten palabras con el catálogo, mientras que la ventaja de Embeddings es interpretar consultas descriptivas, como desventaja, ambos enfoques presentan limitaciones que permiten establecer oportunidades de mejora para trabajos futuros, especialmente en la ampliación de los conjuntos de evaluación y en el fortalecimiento de los mecanismos de selección y validación de las películas recomendadas.